% ═══════════════════════════════════════════════════════════════
% MT-01c: Minitest Artikel + Alphabet
% Spanisch Klasse 7 - Sequenz 1: Empezamos
% ═══════════════════════════════════════════════════════════════
% Dauer: 10 Minuten | Punkte: 26 BE
% ═══════════════════════════════════════════════════════════════

\documentclass[11pt, a4paper]{article}

\usepackage[utf8]{inputenc}
\usepackage[T1]{fontenc}
\usepackage[ngerman]{babel}
\usepackage{helvet}
\renewcommand{\familydefault}{\sfdefault}

\usepackage[margin=2cm]{geometry}
\usepackage{enumitem}
\usepackage{tabularx}
\usepackage{booktabs}
\usepackage{xcolor}
\usepackage{fancyhdr}
\usepackage{lastpage}
\usepackage{amssymb}

% FARBEN
\definecolor{spanisch}{HTML}{c41e3a}
\definecolor{grau}{HTML}{888888}
\definecolor{hellgrau}{HTML}{f5f5f5}

\pagestyle{fancy}
\fancyhf{}
\renewcommand{\headrulewidth}{0.4pt}
\renewcommand{\footrulewidth}{0.4pt}

\lhead{Spanisch -- Klasse 7}
\chead{\textbf{MT-01c}}
\rhead{\today}

\lfoot{10 Minuten}
\cfoot{}
\rfoot{Seite \thepage\ von \pageref{LastPage}}

\newcommand{\punktebox}[1]{\hfill\fbox{\small \_\_\_/#1}}
\newcommand{\luecke}[1]{\rule{#1}{0.4pt}}
\newcommand{\es}[1]{\textcolor{spanisch}{\textbf{#1}}}

\begin{document}

% ═══════════════════════════════════════════════════════════════
% KOPFBEREICH
% ═══════════════════════════════════════════════════════════════

\begin{center}
    {\Large\bfseries\textcolor{spanisch}{Minitest: Artikel + Alphabet}}
\end{center}

\vspace{5pt}

\noindent
\begin{tabularx}{\textwidth}{|X|X|X|}
\hline
\textbf{Name:} \luecke{4cm} &
\textbf{Klasse:} 7 &
\textbf{Datum:} \luecke{2cm} \\
\hline
\end{tabularx}

\vspace{10pt}

\noindent
\fcolorbox{spanisch}{hellgrau}{%
\parbox{\dimexpr\textwidth-2\fboxsep-2\fboxrule}{%
\textbf{Zeit:} 10 Minuten | \textbf{Punkte:} 26 BE | \textbf{Hilfsmittel:} keine
}}

\vspace{15pt}

% ═══════════════════════════════════════════════════════════════
% AUFGABE 1: Singular-Artikel (6 BE)
% ═══════════════════════════════════════════════════════════════

\noindent\textbf{\textcolor{spanisch}{Aufgabe 1 -- Bestimmte Artikel (Singular)}} \punktebox{6}

\noindent Schreibe den richtigen Artikel: \es{el} oder \es{la}.

\vspace{10pt}

\begin{tabularx}{\textwidth}{|l|X|l|X|}
\hline
\luecke{1.5cm} & libro (Buch) & \luecke{1.5cm} & mesa (Tisch) \\[0.5em]
\hline
\luecke{1.5cm} & cuaderno (Heft) & \luecke{1.5cm} & silla (Stuhl) \\[0.5em]
\hline
\luecke{1.5cm} & chico (Junge) & \luecke{1.5cm} & chica (Mädchen) \\[0.5em]
\hline
\end{tabularx}

\vspace{15pt}

% ═══════════════════════════════════════════════════════════════
% AUFGABE 2: Plural-Artikel (6 BE)
% ═══════════════════════════════════════════════════════════════

\noindent\textbf{\textcolor{spanisch}{Aufgabe 2 -- Bestimmte Artikel (Plural)}} \punktebox{6}

\noindent Schreibe den richtigen Plural-Artikel: \es{los} oder \es{las}.

\vspace{10pt}

\begin{tabularx}{\textwidth}{|l|X|l|X|}
\hline
\luecke{1.5cm} & libros (Bücher) & \luecke{1.5cm} & mesas (Tische) \\[0.5em]
\hline
\luecke{1.5cm} & chicos (Jungen) & \luecke{1.5cm} & sillas (Stühle) \\[0.5em]
\hline
\luecke{1.5cm} & profesores (Lehrer) & \luecke{1.5cm} & profesoras (Lehrerinnen) \\[0.5em]
\hline
\end{tabularx}

\vspace{15pt}

% ═══════════════════════════════════════════════════════════════
% AUFGABE 3: Genus-Regel (4 BE)
% ═══════════════════════════════════════════════════════════════

\noindent\textbf{\textcolor{spanisch}{Aufgabe 3 -- Genus-Regeln}} \punktebox{4}

\noindent Ergänze die Regeln. Kreuze an: \textbf{maskulin (el)} oder \textbf{feminin (la)}.

\vspace{10pt}

\noindent
a) Nomen auf \textbf{-o} sind meist:

\hspace{1cm} $\square$ maskulin (el) \hspace{2cm} $\square$ feminin (la)

\vspace{0.8em}

\noindent
b) Nomen auf \textbf{-a} sind meist:

\hspace{1cm} $\square$ maskulin (el) \hspace{2cm} $\square$ feminin (la)

\vspace{15pt}

% ═══════════════════════════════════════════════════════════════
% AUFGABE 4: Alphabet Buchstabennamen (6 BE)
% ═══════════════════════════════════════════════════════════════

\noindent\textbf{\textcolor{spanisch}{Aufgabe 4 -- Spanische Buchstabennamen}} \punktebox{6}

\noindent Wie heißen diese Buchstaben auf Spanisch? Schreibe den Namen.

\vspace{10pt}

\begin{enumerate}[label=\alph*)]
    \item \textbf{H} = \luecke{4cm}
    \item \textbf{J} = \luecke{4cm}
    \item \textbf{Ñ} = \luecke{4cm}
\end{enumerate}

\vspace{15pt}

% ═══════════════════════════════════════════════════════════════
% AUFGABE 5: Buchstabieren (4 BE)
% ═══════════════════════════════════════════════════════════════

\noindent\textbf{\textcolor{spanisch}{Aufgabe 5 -- Buchstabieren}} \punktebox{4}

\noindent Buchstabiere das Wort \es{HOLA} auf Spanisch. Schreibe die Buchstabennamen mit Bindestrichen.

\vspace{10pt}

\noindent
\es{HOLA} = \luecke{10cm}

% ═══════════════════════════════════════════════════════════════
% PUNKTEÜBERSICHT
% ═══════════════════════════════════════════════════════════════

\vspace{20pt}
\hrule
\vspace{10pt}

\noindent
\begin{tabularx}{\textwidth}{|X|c|c|c|c|c||c|}
\hline
\textbf{Aufgabe} & 1 & 2 & 3 & 4 & 5 & \textbf{Gesamt} \\
\hline
\textbf{max. Punkte} & 6 & 6 & 4 & 6 & 4 & \textbf{26} \\
\hline
\textbf{erreicht} & & & & & & \\[0.6em]
\hline
\end{tabularx}

\vspace{10pt}

\begin{flushright}
\textbf{Note:} \luecke{2cm}
\end{flushright}

\end{document}
